\documentclass[a4paper,12pt]{article}
\usepackage{TechIng}

\usepackage[]{graphicx}  
\DeclareGraphicsExtensions{.pdf,.jpeg,.jpg, .png, .eps, .tiff}
\usepackage{subfigure}

%========================================================================================
%%%% ATTENTION !!!
%%% Le package TechIng utilise biblatex avec biber en tant que compilateur !!
\addbibresource{bibreport.bib}


%%%
\fancypagestyle{plain}{\pagestyle{fancy}}
\parindent 0pt

%========================================================================================
% TITRE DU DOCUMENT
%========================================================================================
%%%%% Définitions du titre du document; à renseigner
\def\TIdocumentFrenchTitle{Titre principal}
\def\TIdocumentEnglishTitle{Title (traduction du titre et sous-titre en anglais)}
\def\TIdocumentSubTitle{Sous-titre (optionnel)\newline{\small \hl{\rmfamily Remarque : l'ensemble du titre et sous-titre ne doit pas d\'epasser \textbf{70 signes espaces compris}}}}

%%%% Si un sous-titre existe, utiliser cette définition du titre, sinon commenter
\newcommand{\TImaketitle}{%
	\begin{flushleft}
		{\sffamily\LARGE\textcolor{TI@blue}{\textbf{\TIdocumentFrenchTitle}}} \\[2mm]
		{\sffamily\large\textcolor{TI@blue}{\TIdocumentSubTitle}} \\[4mm]
	\end{flushleft}
}

%%%%% Si aucun sous-titre n'existe, utiliser cette définition du titre, sinon commenter
%\newcommand{\TImaketitle}{%
%	\begin{flushleft}
%		{\sffamily\LARGE\textcolor{TI@blue}{\textbf{\TIdocumentFrenchTitle}}} \\[4mm]
%	\end{flushleft}
%}

%========================================================================================
%========================================================================================
\begin{document}
\selectlanguage{french}

%%%%% Debut entête du document

\TImaketitle
\sffamily\Large\it\textcolor{TI@blue}{\TIdocumentEnglishTitle}\normalsize\normalfont \\[10mm]

\newcommand\TImakeAllAuthors{
	%%% Ajouter autant d'auteurs que nécessaire
	%% Auteur 1
	\def\TIdocumentAuthorName{Pr\'enom NOM Auteur 1}
	\def\TIdocumentAuthorJob{Fonction Auteur 1}
	\def\TIdocumentAuthorAffiliation{Affiliation Auteur 1 \hl{(nom du laboratoire, de l'organisme de recherche ou de l'entreprise, ville, pays)}}
	\TImakeAuthor
	
	%% Auteur 2
	\def\TIdocumentAuthorName{Pr\'enom NOM Auteur 2}
	\def\TIdocumentAuthorJob{Fonction Auteur 2}
	\def\TIdocumentAuthorAffiliation{Affiliation Auteur 2 \hl{(nom du laboratoire, de l'organisme de recherche ou de l'entreprise, ville, pays)}}
	\TImakeAuthor
}
\TImakeAllAuthors

%%%%% Résumé en français
\begin{AbstractFrench}
\textbf{(env. 500 - 700 signes) : R\'esum\'e du texte}


\begin{Remark}%%% Commentaires à supprimer dans la version finale
	\hl{le r\'esum\'e doit \^etre concis et factuel. Il doit exposer bri\`evement la probl\'ematique et le contenu de l'article. Le r\'esum\'e doit pouvoir \^etre lu ind\'ependamment de l'article. Pour cette raison, il ne doit pas contenir de r\'ef\'erences bibliographiques, ni de liens vers d'autres articles.\\
	Le \og nous\fg~est \`a \'eviter (ex. : \og Nous \'etudierons dans cet article...\fg) ; il faut privil\'egier le \og il\fg~impersonnel (ex.~: \og Dans cet article, il est question...\fg~ou bien simplement \og Cet article traite...\fg).}
\end{Remark}
\end{AbstractFrench}

%%%%% Résumé en anglais
\begin{AbstractEnglish}
	{(env. 500 - 700 signs) : Text abstract}
\end{AbstractEnglish}

%%%%% Mots-clés en français
\begin{FrenchKeyWords}
	\textbf{(env. 3 - 4)}
	
	Mot-cl\'e 1, Mot-cl\'e 2, Mot-cl\'e 3, Mot-cl\'e 4.
\end{FrenchKeyWords}

%%%%% Mots-clés en anglais
\begin{EnglishKeyWords}
	\textbf{(env. 3 - 4)}
\end{EnglishKeyWords}


\vspace{1cm}
\begin{Remark}%%% Commentaires à supprimer dans la version finale
	\hl{L'article doit comprendre une \textbf{traduction en anglais - du titre de l'article, - du r\'esum\'e, - des diff\'erents mots-clefs}. Ces \'el\'ements sont indispensables pour le r\'ef\'erencement et l'indexation des articles dans les bases de donn\'ees bibliographiques internationales.}
\end{Remark}

%%%%% Fin de l'entête du document

%========================================================================================
%%%%% Table des matières
\clearpage
\tableofcontents

%========================================================================================
%========================================================================================
%%%%% Introduction
\begin{Introduction}
	%%% Commentaires à supprimer dans la version finale
	Tout article doit d\'ebuter par une introduction de 1 500 \`a 5 000 caract\`eres (espaces compris) \textbf{comprenant :}
	\begin{itemize}
	\item \textbf{la d\'efinition du concept ou de la technologie principale d\'ecrite dans l'article, }
	\item \textbf{le contexte, }
	\item \textbf{les probl\'ematiques. }
	\end{itemize}

	Cette introduction est lisible par tout utilisateur du site Internet www.techniques-ingenieur.fr, m\^eme non abonn\'e.\newline
	
	Ce pr\'eambule doit permettre au lecteur de savoir quel est l'objectif de l'article et sous quel angle il est \'ecrit. Il replace le sujet dans l'environnement technico-\'economique. L'introduction rappellera les domaines d'application, les atouts et les limites de la technique trait\'ee.\newline
	
	\textbf{Ne comporte ni appels bibliographiques, ni figures, ni tableaux.}\newline
	
	Nota : le lecteur trouvera en fin d'article un glossaire des termes et expressions importants de l'article, ainsi qu'un tableau des sigles, notations et symboles utilis\'es tout au long de l'article.
\end{Introduction}

%========================================================================================
%%%% Corps du texte
\section{Titre de niveau 1}\label{sec:niveau1}

Texte courant --  Texte courant -- Texte courant -- Texte courant -- Texte courant\newline

%%% Environnement Definitions
\begin{Definition}
	Les \textbf{d\'efinitions} \`a retenir sont mises en exergue avec une trame bleue.\newline
\end{Definition}


Les \textbf{figures} et les \textbf{tableaux} doivent \^etre num\'erot\'es de 1 \`a $n$ et appel\'es dans le texte.\newline

Les \textbf{\'equations} peuvent \^etre num\'erot\'ees de 1 \`a $n$ et appel\'ees dans le texte.\newline

Les \textbf{r\'ef\'erences bibliographiques} num\'erot\'ees de [1] \`a [$n$] (utiliser la commande \verb?\cite? en \LaTeX) concernent les articles (hormis ceux de notre collection), ouvrages, actes de colloques et th\`eses et ne concernent pas les normes, sites, lois, brevets, congr\`es (cf. \og Pour en savoir plus\fg). Les \textbf{articles de notre collection} sont appel\'es en indiquant leur r\'ef\'erence interne, par exemple [D2901] (\`a r\'ef\'erencer manuellement).\newline

Les \textbf{sigles} doivent \^etre d\'evelopp\'es \`a leur premi\`ere occurrence.\newline

%%% Environnement Exemple
\begin{Exemple}
	\label{exemple:ex1}
	\titreexemple{\ref{exemple:ex1} : Titre de l'exemple}
	Les \textbf{exemples} pr\'esent\'es de cette mani\`ere sont les bienvenus.\\
	Pour les \'equations dans les exemples, il est possible d'utiliser les environnements et commandes suivants
	\begin{equationExemple}
		a^2-b^2=(a+b)(a-b)
	\end{equationExemple}
	ou
	\begin{equation*}
		a^2+b^2+2\,a\,b=(a+b)^2 \numberthis
	\end{equation*}
	ou encore
	\begin{align*}
		E=m\,c^2 \numberthis\\ 
		f=m\ddot{x}\, \numberthis
	\end{align*}
\end{Exemple}


%%% Environnement Figures
\begin{figure}[t]
	\centering
	\includegraphics[width=\linewidth]{FigureTest}
	\caption{Titre de figure}
	\label{fig:fig1}
	\begin{flushleft}
		\begin{Remark}%%% Commentaires à supprimer dans la version finale
			\hl{les titres de figure ne doivent comporter ni \'equation ni symbole. Les l\'egendes doivent \^etre en fran\c{c}ais.}
		\end{Remark}
	\end{flushleft}
\end{figure}

%%% Environnement Tableaux
\begin{table}[h]
	\centering
	\caption{Titre de tableau.}
	\begin{tabular}{|c|c|c|}
		\hline
		~~~~~~~~~~~T\^eti\`ere1~~~~~~~~~~~ 	& ~~~~~~~~~~~T\^eti\`ere2~~~~~~~~~~~ 	& ~~~~~~~~~~~T\^eti\`ere3~~~~~~~~~~~\\
		\hline
		123 		& 456		& 789\\
		\hline
	\end{tabular}
	\label{tab:tab1}
	\begin{flushleft}
		\begin{Remark}%%% Commentaires à supprimer dans la version finale
			\begin{itemize}
					\item \hl{Les titres de tableau ne doivent comporter ni \'equation ni symbole.}
					\item \hl{Les unit\'es doivent figurer uniquement en t\^ete de colonne et de ligne.}
			\end{itemize}
		\end{Remark}
	\end{flushleft}
\end{table}

%%% A supprimer dans la version finale
\fcolorbox{black}{TI@paleyellow}{
	\begin{minipage}{\linewidth}
		\begin{center}
			\hl{\textbf{Photos, figures et tableaux}}
		\end{center}
		\begin{itemize}
			\item Toutes les photos, toutes les figures et tous les tableaux doivent \^etre de pr\'ef\'erence de votre composition, les figures si possible au format PowerPoint.
			\item Toutes les photos, toutes les figures et tous les tableaux puis\'es ailleurs feront l'objet de votre part d'une demande d'autorisation de reproduction au fur et \`a mesure de la progression de la r\'edaction de votre article. Indiquer le copyright pour chacun. Refaire les figures si besoin.
			\item Les photos devront \^etre fournies en fichier s\'epar\'e au format jpeg en r\'esolution 300 dpi minimum (haute qualit\'e).
			\item Les copies d'\'ecran sont \`a \'eviter.
		\end{itemize}
	\end{minipage}
}\newline

%%% Environnement equations
\begin{equation}
	PV=nRT
	\label{eq:eq1}
\end{equation}
\begin{Remark}%%% Commentaires à supprimer dans la version finale
	\hl{ les formules et \'equations ne doivent pas comporter de titre.}\newline\newline
\end{Remark}


%%% Environnement Cadre
\begin{Cadre}
	\begin{center}
		\textbf{Titre de l'encadr\'e (facultatif)}
	\end{center}
	Les informations \`a sortir du texte courant sont plac\'ees dans un encadr\'e.
\end{Cadre}

%%% Environnement Nota
\hl{\textbf{Les notes de bas de page sont proscrites.}}%%% Commentaires à supprimer dans la version finale

\begin{Nota}
	ce nota se substitue aux notes de bas de page.
\end{Nota}


\subsection{Titre de niveau 2}\label{sec:niveau2}
Texte courant --  Texte courant -- Texte courant -- Texte courant -- Texte courant

\subsubsection{Titre de niveau 3}\label{sec:niveau3}
Texte courant --  Texte courant -- Texte courant -- Texte courant -- Texte courant

\paragraph{Titre de niveau 4}\label{sec:niveau4}~\par
Texte courant --  Texte courant -- Texte courant -- Texte courant -- Texte courant\newline



%%%% Conclusion
\section*{Conclusion}
Dans sa conclusion, l'auteur ouvrira sur l'avenir et les \'evolutions technologiques possibles, mais il \'evitera de donner des \'ech\'eances pr\'ecises afin de ne pas dater l'article.



%%%% Glossaire
\section*{Glossaire}
L'auteur d\'efinira les termes ou expressions les plus importants de l'article. La traduction des termes en anglais est conseill\'ee.\newline

\hl{Exemples : }\newline

%%% Environnement Glossaire
\begin{Glossaire}{thermogravim\'etrie}{thermogravimetry}
	Technique permettant de mesurer, en fonction du temps ou de la temp\'erature, la masse d'un \'echantillon chauff\'e ou refroidi selon un programme de temp\'erature d\'etermin\'e.
\end{Glossaire}

\begin{Acronyme}{ICP-MS}{inductively coupled plasma-mass spectrometer}
	Spectrom\`etre de masse utilis\'e pour d\'eterminer la nature d'un mat\'eriau, qui analyse les ions qui se forment lorsque les \'el\'ements d'un \'echantillon traversent un plasma \`a couplage inductif.
\end{Acronyme}



%%%% Sigles, notations et symboles
\section*{Sigles, notations et symboles}
\hl{Exemples : }\newline

\begin{table}[h]
	\centering
	\begin{tabular}{|p{\dimexpr 0.2\linewidth-2\tabcolsep}|p{\dimexpr 0.6\linewidth-2\tabcolsep}|p{\dimexpr 0.2\linewidth-2\tabcolsep}|}
		\hline
		\textbf{Symbole} & \textbf{Description} & \textbf{Unit\'e}\\
		\hline
		A 	& Section d'\'ecoulement de la soupape de s\^uret\'e		& mm$^2$\\
		\hline
		C 	& Fonction du coefficient isentropique		& -- \\
		\hline
		$K_b$ 	& Facteur de correction du d\'ebit th\'eorique pour un \'ecoulement sub-critique		& -- \\
		\hline
		$K_d$ 	& Coefficient de d\'ebit certifi\'e		& -- \\
		\hline
		$K_{dr}$ 	& Coefficient de d\'ebit apr\`es abattement ($K_d \times 0.9$)		& -- \\
		\hline
		$\ldots$ 	& $\ldots$		& -- \\
		\hline
	\end{tabular}
	\label{tab:tabsignes}
\end{table}

\clearpage




%========================================================================================
%========================================================================================
%%%%% Début pied du document

\subsection*{Pour en savoir plus\newline}

%%%% Ne pas supprimer ces trois commandes
\TImaketitle
\TImakeAllAuthors
\normalfont

%========================================================================================
%	BIBLIOGRAPHY
%========================================================================================

%% Seuls les articles, ouvrages scientifiques et chapitres d'ouvrages, et brevets, sont à mettre dans le fichier de références bibliographiques "bibreport.bib"
%% Le reste (outils logiciels, Sites Internet, Evénements, Normes et standards, Réglementation, etc.) sont à lister dans les sections adhoc ci-dessous

\nocite{Curriebook,patent_test}

\defbibfilter{nopatent}{
	not type=patent
}
\printbibliography[filter=nopatent,title=Sources bibliographiques]
\hl{(ouvrages, articles, actes de colloques, th\`eses et m\'emoires, appel\'es dans le texte -- donc num\'erot\'ees -- ou non)}

\subsubsection*{\`A lire \'egalement dans nos bases}\rmfamily

\hl{(articles des Techniques de l'Ing\'enieur. R\'ef\'erences non num\'erot\'ees, qui peuvent au besoin \^etre appel\'ees dans le texte par leur num\'ero de r\'ef\'erence entre crochets, par exemple [D2901])}

MOREL Florent. -- Commande pr\'edictive des machines \'electriques tournantes. [D2901], Convertisseurs et machines \'electriques (2009).

\section*{Outils logiciels}

\hl{(logiciels)}

CISIA juin 2000 \textit{Le Bay\'esien} (version pour Windows Vista), [Logiciel] CISIA 1 avenue Herbillon, 94160 Saint-Mand\'e, France\\
Autre logiciel

\section*{Sites Internet}
\hl{(portails, pages personnelles, forums, listes de discussion)}

AMIRE Gilles Un Nouveau Guide Internet \url{http://www.imaginet.fr/ime/toc.htm} (page consult\'ee le 23 juin 1999)\\
Autre site web

\section*{\'Ev\'enements}
\hl{(congr\`es, salons, colloques, journ\'ees d'\'etude)}

Salon : Batimat Reed Expositions Salon international de la construction, a lieu tous les deux ans \`a Paris (ann\'ees impaires). \url{http://www.batimat.com}\\
Autre \'ev\'enement

\section*{Normes et standards}
ISO \textit{Acoustique : m\'ethode de calcul du niveau d'isotonie} ISO 532 1975 ISO, 1975.

AFNOR \textit{Tuiles en terre cuite pour pos\'e en discontinu. D\'etermination des caract\'eristiques physiques. Partie 1 : essai d'imperm\'eabilit\'e} NF EN 539-1 AFNOR, 1994



\section*{R\'eglementation}
Loi No 83-634 du 13 juillet 1983 portant droits et obligations... (version consolid\'ee du 3 mars 2002).

D\'ecret No 2000-44 du 13 janvier 2000 portant... (version consolid\'ee au 5 octobre 2007) JO No 11 du 14 janvier 2000 page 369 NOR : FPPA9910013D


% Pour les brevets, on les distinguera du reste des publications en les référençant avec un B devant
% Pour réussir à faire cela, il faut rajouter dans les keyword des références associées aux fichier bibtex la lettre B
% Exemple du brevet référencé dans le fichier bibreport.bib fourni au template
% @Patent{patent_test,
%	nationality = {CA},
%	number      = {CA 2582058},
%	title       = {Proc\'ed\'e d'oxydation \'electrolytique du c\'erium et ensemble d'\'electrolyse pour sa mise en oeuvre},
%	keywords    = {B},
% }
\newrefcontext[labelprefix=B]
\printbibliography[type=patent,resetnumbers=true,title=Brevets,keyword=B]





\section*{Annuaire}
\subsubsection*{Constructeurs -- Fournisseurs -- Distributeurs (liste non exhaustive)}
\hl{Avec adresse Internet}

Presses \`a injecter horizontales \'electriques

Arburg SAS Hawa\"i distribu\'e par Arburg France \url{http://www.arburg.fr}

\subsubsection*{Organismes -- F\'ed\'erations -- Associations (liste non exhaustive)}
\hl{Avec adresse Internet}

Organisme professionnel de pr\'evention du b\^atiment et des travaux publics OPPBT \url{http://www.oppbt.fr}


\subsubsection*{Documentation -- Formation -- S\'eminaires (liste non exhaustive)}
\hl{Avec adresse Internet}

S\'eminaire MCBT. \textit{Films supraconducteurs \`a HTc et THTc : r\^ole des interfaces dans leur \'elaboration par voie chimique.} Philippe Odier mardi 20 novembre 2007 Institut N\'eel \url{http://neel.cnrs.fr/spip.php?article1324}

\subsubsection*{Laboratoires -- Bureaux d'\'etudes -- \'Ecoles -- Centres de recherche  (liste non exhaustive)}
\hl{Avec adresse Internet}

Centre technique et scientifique du b\^atiment CSTB \url{http//www.cstb.fr}

\section*{Donn\'ees statistiques et \'economiques}
\hl{Texte libre, pouvant contenir des tableaux, figures (dont cartes et camemberts)}



%----------------------------------------------------------------------------------------

\label{LastPage}

\end{document}